% Options for packages loaded elsewhere
\PassOptionsToPackage{unicode}{hyperref}
\PassOptionsToPackage{hyphens}{url}
%
\documentclass[
]{article}
\usepackage{amsmath,amssymb}
\usepackage{iftex}
\ifPDFTeX
  \usepackage[T1]{fontenc}
  \usepackage[utf8]{inputenc}
  \usepackage{textcomp} % provide euro and other symbols
\else % if luatex or xetex
  \usepackage{unicode-math} % this also loads fontspec
  \defaultfontfeatures{Scale=MatchLowercase}
  \defaultfontfeatures[\rmfamily]{Ligatures=TeX,Scale=1}
\fi
\usepackage{lmodern}
\ifPDFTeX\else
  % xetex/luatex font selection
\fi
% Use upquote if available, for straight quotes in verbatim environments
\IfFileExists{upquote.sty}{\usepackage{upquote}}{}
\IfFileExists{microtype.sty}{% use microtype if available
  \usepackage[]{microtype}
  \UseMicrotypeSet[protrusion]{basicmath} % disable protrusion for tt fonts
}{}
\makeatletter
\@ifundefined{KOMAClassName}{% if non-KOMA class
  \IfFileExists{parskip.sty}{%
    \usepackage{parskip}
  }{% else
    \setlength{\parindent}{0pt}
    \setlength{\parskip}{6pt plus 2pt minus 1pt}}
}{% if KOMA class
  \KOMAoptions{parskip=half}}
\makeatother
\usepackage{xcolor}
\usepackage[margin=1in]{geometry}
\usepackage{color}
\usepackage{fancyvrb}
\newcommand{\VerbBar}{|}
\newcommand{\VERB}{\Verb[commandchars=\\\{\}]}
\DefineVerbatimEnvironment{Highlighting}{Verbatim}{commandchars=\\\{\}}
% Add ',fontsize=\small' for more characters per line
\usepackage{framed}
\definecolor{shadecolor}{RGB}{248,248,248}
\newenvironment{Shaded}{\begin{snugshade}}{\end{snugshade}}
\newcommand{\AlertTok}[1]{\textcolor[rgb]{0.94,0.16,0.16}{#1}}
\newcommand{\AnnotationTok}[1]{\textcolor[rgb]{0.56,0.35,0.01}{\textbf{\textit{#1}}}}
\newcommand{\AttributeTok}[1]{\textcolor[rgb]{0.13,0.29,0.53}{#1}}
\newcommand{\BaseNTok}[1]{\textcolor[rgb]{0.00,0.00,0.81}{#1}}
\newcommand{\BuiltInTok}[1]{#1}
\newcommand{\CharTok}[1]{\textcolor[rgb]{0.31,0.60,0.02}{#1}}
\newcommand{\CommentTok}[1]{\textcolor[rgb]{0.56,0.35,0.01}{\textit{#1}}}
\newcommand{\CommentVarTok}[1]{\textcolor[rgb]{0.56,0.35,0.01}{\textbf{\textit{#1}}}}
\newcommand{\ConstantTok}[1]{\textcolor[rgb]{0.56,0.35,0.01}{#1}}
\newcommand{\ControlFlowTok}[1]{\textcolor[rgb]{0.13,0.29,0.53}{\textbf{#1}}}
\newcommand{\DataTypeTok}[1]{\textcolor[rgb]{0.13,0.29,0.53}{#1}}
\newcommand{\DecValTok}[1]{\textcolor[rgb]{0.00,0.00,0.81}{#1}}
\newcommand{\DocumentationTok}[1]{\textcolor[rgb]{0.56,0.35,0.01}{\textbf{\textit{#1}}}}
\newcommand{\ErrorTok}[1]{\textcolor[rgb]{0.64,0.00,0.00}{\textbf{#1}}}
\newcommand{\ExtensionTok}[1]{#1}
\newcommand{\FloatTok}[1]{\textcolor[rgb]{0.00,0.00,0.81}{#1}}
\newcommand{\FunctionTok}[1]{\textcolor[rgb]{0.13,0.29,0.53}{\textbf{#1}}}
\newcommand{\ImportTok}[1]{#1}
\newcommand{\InformationTok}[1]{\textcolor[rgb]{0.56,0.35,0.01}{\textbf{\textit{#1}}}}
\newcommand{\KeywordTok}[1]{\textcolor[rgb]{0.13,0.29,0.53}{\textbf{#1}}}
\newcommand{\NormalTok}[1]{#1}
\newcommand{\OperatorTok}[1]{\textcolor[rgb]{0.81,0.36,0.00}{\textbf{#1}}}
\newcommand{\OtherTok}[1]{\textcolor[rgb]{0.56,0.35,0.01}{#1}}
\newcommand{\PreprocessorTok}[1]{\textcolor[rgb]{0.56,0.35,0.01}{\textit{#1}}}
\newcommand{\RegionMarkerTok}[1]{#1}
\newcommand{\SpecialCharTok}[1]{\textcolor[rgb]{0.81,0.36,0.00}{\textbf{#1}}}
\newcommand{\SpecialStringTok}[1]{\textcolor[rgb]{0.31,0.60,0.02}{#1}}
\newcommand{\StringTok}[1]{\textcolor[rgb]{0.31,0.60,0.02}{#1}}
\newcommand{\VariableTok}[1]{\textcolor[rgb]{0.00,0.00,0.00}{#1}}
\newcommand{\VerbatimStringTok}[1]{\textcolor[rgb]{0.31,0.60,0.02}{#1}}
\newcommand{\WarningTok}[1]{\textcolor[rgb]{0.56,0.35,0.01}{\textbf{\textit{#1}}}}
\usepackage{graphicx}
\makeatletter
\newsavebox\pandoc@box
\newcommand*\pandocbounded[1]{% scales image to fit in text height/width
  \sbox\pandoc@box{#1}%
  \Gscale@div\@tempa{\textheight}{\dimexpr\ht\pandoc@box+\dp\pandoc@box\relax}%
  \Gscale@div\@tempb{\linewidth}{\wd\pandoc@box}%
  \ifdim\@tempb\p@<\@tempa\p@\let\@tempa\@tempb\fi% select the smaller of both
  \ifdim\@tempa\p@<\p@\scalebox{\@tempa}{\usebox\pandoc@box}%
  \else\usebox{\pandoc@box}%
  \fi%
}
% Set default figure placement to htbp
\def\fps@figure{htbp}
\makeatother
\setlength{\emergencystretch}{3em} % prevent overfull lines
\providecommand{\tightlist}{%
  \setlength{\itemsep}{0pt}\setlength{\parskip}{0pt}}
\setcounter{secnumdepth}{-\maxdimen} % remove section numbering
\usepackage{booktabs}
\usepackage{longtable}
\usepackage{array}
\usepackage{multirow}
\usepackage{wrapfig}
\usepackage{float}
\usepackage{colortbl}
\usepackage{pdflscape}
\usepackage{tabu}
\usepackage{threeparttable}
\usepackage{threeparttablex}
\usepackage[normalem]{ulem}
\usepackage{makecell}
\usepackage{xcolor}
\usepackage{bookmark}
\IfFileExists{xurl.sty}{\usepackage{xurl}}{} % add URL line breaks if available
\urlstyle{same}
\hypersetup{
  pdftitle={Workflow RNA-seq Analysis},
  pdfauthor={Mariana Mescouto Lopes},
  hidelinks,
  pdfcreator={LaTeX via pandoc}}

\title{Workflow RNA-seq Analysis}
\usepackage{etoolbox}
\makeatletter
\providecommand{\subtitle}[1]{% add subtitle to \maketitle
  \apptocmd{\@title}{\par {\large #1 \par}}{}{}
}
\makeatother
\subtitle{Analysis of count matrix data}
\author{Mariana Mescouto Lopes}
\date{2026}

\begin{document}
\maketitle

{
\setcounter{tocdepth}{2}
\tableofcontents
}
\subsection{Overview}\label{overview}

This workflow demonstrates a basic RNA-seq analysis pipeline.\\
To run the analysis, update the file paths according to your local
environment.

This workflow performs: - Data preprocessing - Differential expression
analysis - Data visualization - Functional enrichment analysis (GSEA)

\section{1. Packges}\label{packges}

To install packages, use: install.packages(``package\_name'')

Then load them using: library(package\_name)

\begin{Shaded}
\begin{Highlighting}[]
\FunctionTok{suppressPackageStartupMessages}\NormalTok{(\{}
  \CommentTok{\# Diferential analysis}
  \FunctionTok{library}\NormalTok{(DESeq2)}
  \FunctionTok{library}\NormalTok{(edgeR) }
  \FunctionTok{library}\NormalTok{(limma)}
  \FunctionTok{library}\NormalTok{(fdrtool)}
  
  \CommentTok{\# Data manipulation}
  \FunctionTok{library}\NormalTok{(dplyr)}
  \FunctionTok{library}\NormalTok{(readr)}
  \FunctionTok{library}\NormalTok{(readxl)}
  \FunctionTok{library}\NormalTok{(tibble)}
  
  \CommentTok{\# Visualization}
  \FunctionTok{library}\NormalTok{(ggplot2)}
  \FunctionTok{library}\NormalTok{(pheatmap)}
  \FunctionTok{library}\NormalTok{(ggrepel)}
  \FunctionTok{library}\NormalTok{(RColorBrewer)}
  \FunctionTok{library}\NormalTok{(corrplot)}
  \FunctionTok{library}\NormalTok{(plotly)}
  
  \CommentTok{\# Functional analysis}
  \FunctionTok{library}\NormalTok{(fgsea)}
  \FunctionTok{library}\NormalTok{(biomaRt)}
  
  \CommentTok{\# Tables}
  \FunctionTok{library}\NormalTok{(DT)}
  \FunctionTok{library}\NormalTok{(kableExtra)}
\NormalTok{\})}
\end{Highlighting}
\end{Shaded}

\section{2. Data preprocessing}\label{data-preprocessing}

Load count matrix and sample metadata

\begin{Shaded}
\begin{Highlighting}[]
\NormalTok{df }\OtherTok{\textless{}{-}} \FunctionTok{read\_delim}\NormalTok{(}\StringTok{"rnaseq\_data.csv"}\NormalTok{, }\AttributeTok{delim =} \StringTok{";"}\NormalTok{, }\AttributeTok{show\_col\_types =} \ConstantTok{FALSE}\NormalTok{)}
\NormalTok{metadata }\OtherTok{\textless{}{-}} \FunctionTok{read\_excel}\NormalTok{(}\StringTok{"samples.xlsx"}\NormalTok{)}

\CommentTok{\# Dataset information}
\FunctionTok{cat}\NormalTok{(}\FunctionTok{sprintf}\NormalTok{(}\StringTok{"Raw dataset: \%d genes × \%d columns}\SpecialCharTok{\textbackslash{}n}\StringTok{"}\NormalTok{, }\FunctionTok{nrow}\NormalTok{(df), }\FunctionTok{ncol}\NormalTok{(df)))}
\end{Highlighting}
\end{Shaded}

\begin{verbatim}
## Raw dataset: 35670 genes × 19 columns
\end{verbatim}

\begin{Shaded}
\begin{Highlighting}[]
\FunctionTok{cat}\NormalTok{(}\FunctionTok{sprintf}\NormalTok{(}\StringTok{"Metadata: \%d samples}\SpecialCharTok{\textbackslash{}n}\StringTok{"}\NormalTok{, }\FunctionTok{nrow}\NormalTok{(metadata)))}
\end{Highlighting}
\end{Shaded}

\begin{verbatim}
## Metadata: 18 samples
\end{verbatim}

\subsection{Sample selection and gene
formatting}\label{sample-selection-and-gene-formatting}

\begin{Shaded}
\begin{Highlighting}[]
\NormalTok{ctrl\_samples }\OtherTok{\textless{}{-}} \FunctionTok{grep}\NormalTok{(}\StringTok{"CON\_"}\NormalTok{, }\FunctionTok{colnames}\NormalTok{(df), }\AttributeTok{value =} \ConstantTok{TRUE}\NormalTok{)}
\NormalTok{treat\_samples }\OtherTok{\textless{}{-}} \FunctionTok{grep}\NormalTok{(}\StringTok{"BET\_"}\NormalTok{, }\FunctionTok{colnames}\NormalTok{(df), }\AttributeTok{value =} \ConstantTok{TRUE}\NormalTok{)}

\FunctionTok{cat}\NormalTok{(}\FunctionTok{sprintf}\NormalTok{(}\StringTok{"Samples: \%d control + \%d treated}\SpecialCharTok{\textbackslash{}n}\StringTok{"}\NormalTok{, }
            \FunctionTok{length}\NormalTok{(ctrl\_samples), }\FunctionTok{length}\NormalTok{(treat\_samples)))}
\end{Highlighting}
\end{Shaded}

\begin{verbatim}
## Samples: 9 control + 9 treated
\end{verbatim}

\begin{Shaded}
\begin{Highlighting}[]
\CommentTok{\# Set gene IDs as row names}
\NormalTok{genes }\OtherTok{\textless{}{-}}\NormalTok{ df}\SpecialCharTok{$}\NormalTok{gene}
\NormalTok{df }\OtherTok{\textless{}{-}} \FunctionTok{as.matrix}\NormalTok{(df[, }\FunctionTok{c}\NormalTok{(ctrl\_samples, treat\_samples)])}
\FunctionTok{rownames}\NormalTok{(df) }\OtherTok{\textless{}{-}}\NormalTok{ genes}
\end{Highlighting}
\end{Shaded}

\subsection{Filtering low expressed
genes}\label{filtering-low-expressed-genes}

Criteria: keep genes with ≥5 counts in at least 6 samples

\begin{Shaded}
\begin{Highlighting}[]
\NormalTok{filter\_mask }\OtherTok{\textless{}{-}} \FunctionTok{rowSums}\NormalTok{(df }\SpecialCharTok{\textgreater{}=} \DecValTok{5}\NormalTok{) }\SpecialCharTok{\textgreater{}=} \DecValTok{6}
\NormalTok{expr\_genes }\OtherTok{\textless{}{-}}\NormalTok{ genes[filter\_mask]}
\NormalTok{df\_filtered }\OtherTok{\textless{}{-}}\NormalTok{ df[filter\_mask, ]}

\FunctionTok{cat}\NormalTok{(}\FunctionTok{sprintf}\NormalTok{(}\StringTok{"Initial number of genes: \%d}\SpecialCharTok{\textbackslash{}n}\StringTok{"}\NormalTok{, }\FunctionTok{nrow}\NormalTok{(df)))}
\end{Highlighting}
\end{Shaded}

\begin{verbatim}
## Initial number of genes: 35670
\end{verbatim}

\begin{Shaded}
\begin{Highlighting}[]
\FunctionTok{cat}\NormalTok{(}\FunctionTok{sprintf}\NormalTok{(}\StringTok{"Filtered dataset: \%d genes (\%.1f\%\% retained)}\SpecialCharTok{\textbackslash{}n}\StringTok{"}\NormalTok{, }
            \FunctionTok{nrow}\NormalTok{(df\_filtered), }
            \DecValTok{100} \SpecialCharTok{*} \FunctionTok{nrow}\NormalTok{(df\_filtered) }\SpecialCharTok{/} \FunctionTok{nrow}\NormalTok{(df)))}
\end{Highlighting}
\end{Shaded}

\begin{verbatim}
## Filtered dataset: 5764 genes (16.2% retained)
\end{verbatim}

\subsection{Metadata alignment}\label{metadata-alignment}

\begin{Shaded}
\begin{Highlighting}[]
\NormalTok{metadata\_filtered }\OtherTok{\textless{}{-}}\NormalTok{ metadata }\SpecialCharTok{\%\textgreater{}\%}
  \FunctionTok{filter}\NormalTok{(ID }\SpecialCharTok{\%in\%} \FunctionTok{c}\NormalTok{(ctrl\_samples, treat\_samples)) }\SpecialCharTok{\%\textgreater{}\%}
  \FunctionTok{arrange}\NormalTok{(}\FunctionTok{match}\NormalTok{(ID, }\FunctionTok{c}\NormalTok{(ctrl\_samples, treat\_samples)))}
\end{Highlighting}
\end{Shaded}

Save filtered data

\begin{Shaded}
\begin{Highlighting}[]
\FunctionTok{write\_csv}\NormalTok{(}\FunctionTok{as.data.frame}\NormalTok{(df\_filtered), }\StringTok{"results/tables/filtered\_data.csv"}\NormalTok{)}
\end{Highlighting}
\end{Shaded}

\section{3. Exploratory Data Analysis}\label{exploratory-data-analysis}

\subsection{Data distribution}\label{data-distribution}

Evaluate data quality and filtering criteria

\begin{Shaded}
\begin{Highlighting}[]
\NormalTok{total\_expression }\OtherTok{\textless{}{-}} \FunctionTok{rowSums}\NormalTok{(df\_filtered)}

\NormalTok{p\_dist }\OtherTok{\textless{}{-}} \FunctionTok{ggplot}\NormalTok{(}\FunctionTok{data.frame}\NormalTok{(total\_expression), }
                 \FunctionTok{aes}\NormalTok{(}\AttributeTok{x =} \FunctionTok{log10}\NormalTok{(total\_expression }\SpecialCharTok{+} \DecValTok{1}\NormalTok{))) }\SpecialCharTok{+}
  \FunctionTok{geom\_histogram}\NormalTok{(}\AttributeTok{bins =} \DecValTok{50}\NormalTok{, }\AttributeTok{fill =} \StringTok{"steelblue"}\NormalTok{, }\AttributeTok{alpha =} \FloatTok{0.7}\NormalTok{) }\SpecialCharTok{+}
  \FunctionTok{geom\_vline}\NormalTok{(}\AttributeTok{xintercept =} \FunctionTok{log10}\NormalTok{(}\FunctionTok{c}\NormalTok{(}\DecValTok{10}\NormalTok{, }\DecValTok{100}\NormalTok{, }\DecValTok{1000}\NormalTok{)), }
             \AttributeTok{linetype =} \StringTok{"dashed"}\NormalTok{, }\AttributeTok{color =} \StringTok{"red"}\NormalTok{, }\AttributeTok{alpha =} \FloatTok{0.7}\NormalTok{) }\SpecialCharTok{+}
  \FunctionTok{labs}\NormalTok{(}
    \AttributeTok{title =} \StringTok{"Expression distribution"}\NormalTok{,}
    \AttributeTok{subtitle =} \StringTok{"Red lines: 10, 100, 1000 reads"}\NormalTok{,}
    \AttributeTok{x =} \StringTok{"log10(total expression)"}\NormalTok{,}
    \AttributeTok{y =} \StringTok{"Number of genes"}
\NormalTok{  ) }\SpecialCharTok{+}
  \FunctionTok{theme\_minimal}\NormalTok{()}

\FunctionTok{print}\NormalTok{(p\_dist)}
\end{Highlighting}
\end{Shaded}

\pandocbounded{\includegraphics[keepaspectratio]{RNA-seq-pipeline_files/figure-latex/basic-observations-1.pdf}}

Most genes show low expression levels, which is expected in RNA-seq
data. The filtering step removes lowly expressed genes, reducing noise
and improving downstream analysis.

\subsection{PCA}\label{pca}

Explore sample variability and identify potential clustering patterns

\begin{enumerate}
\def\labelenumi{\arabic{enumi})}
\tightlist
\item
  Normalization
\end{enumerate}

\begin{Shaded}
\begin{Highlighting}[]
\NormalTok{df\_log }\OtherTok{\textless{}{-}} \FunctionTok{log2}\NormalTok{(df\_filtered }\SpecialCharTok{+} \DecValTok{1}\NormalTok{)}
\end{Highlighting}
\end{Shaded}

\begin{enumerate}
\def\labelenumi{\arabic{enumi})}
\setcounter{enumi}{1}
\tightlist
\item
  PCA computation
\end{enumerate}

\begin{Shaded}
\begin{Highlighting}[]
\NormalTok{pca }\OtherTok{\textless{}{-}} \FunctionTok{prcomp}\NormalTok{(}\FunctionTok{t}\NormalTok{(df\_log), }\AttributeTok{scale. =} \ConstantTok{TRUE}\NormalTok{)}
\NormalTok{pca\_data }\OtherTok{\textless{}{-}} \FunctionTok{data.frame}\NormalTok{(pca}\SpecialCharTok{$}\NormalTok{x)}

\NormalTok{pca\_data}\SpecialCharTok{$}\NormalTok{sample }\OtherTok{\textless{}{-}} \FunctionTok{rownames}\NormalTok{(pca\_data)}
\NormalTok{pca\_data}\SpecialCharTok{$}\NormalTok{group }\OtherTok{\textless{}{-}} \FunctionTok{ifelse}\NormalTok{(}\FunctionTok{grepl}\NormalTok{(}\StringTok{"CON\_"}\NormalTok{, pca\_data}\SpecialCharTok{$}\NormalTok{sample), }\StringTok{"Control"}\NormalTok{, }\StringTok{"Treated"}\NormalTok{)}
\end{Highlighting}
\end{Shaded}

\begin{enumerate}
\def\labelenumi{\arabic{enumi})}
\setcounter{enumi}{2}
\tightlist
\item
  PCA plot
\end{enumerate}

\begin{Shaded}
\begin{Highlighting}[]
\NormalTok{p\_pca }\OtherTok{\textless{}{-}} \FunctionTok{ggplot}\NormalTok{(pca\_data, }\FunctionTok{aes}\NormalTok{(PC1, PC2, }\AttributeTok{color =}\NormalTok{ group)) }\SpecialCharTok{+}
  \FunctionTok{geom\_point}\NormalTok{(}\AttributeTok{size =} \DecValTok{3}\NormalTok{, }\AttributeTok{alpha =} \FloatTok{0.8}\NormalTok{) }\SpecialCharTok{+}
  \FunctionTok{labs}\NormalTok{(}
    \AttributeTok{title =} \StringTok{"PCA of RNA{-}seq samples"}\NormalTok{,}
    \AttributeTok{x =} \FunctionTok{paste0}\NormalTok{(}\StringTok{"PC1 ("}\NormalTok{, }\FunctionTok{round}\NormalTok{(}\FunctionTok{summary}\NormalTok{(pca)}\SpecialCharTok{$}\NormalTok{importance[}\DecValTok{2}\NormalTok{,}\DecValTok{1}\NormalTok{]}\SpecialCharTok{*}\DecValTok{100}\NormalTok{, }\DecValTok{1}\NormalTok{), }\StringTok{"\%)"}\NormalTok{),}
    \AttributeTok{y =} \FunctionTok{paste0}\NormalTok{(}\StringTok{"PC2 ("}\NormalTok{, }\FunctionTok{round}\NormalTok{(}\FunctionTok{summary}\NormalTok{(pca)}\SpecialCharTok{$}\NormalTok{importance[}\DecValTok{2}\NormalTok{,}\DecValTok{2}\NormalTok{]}\SpecialCharTok{*}\DecValTok{100}\NormalTok{, }\DecValTok{1}\NormalTok{), }\StringTok{"\%)"}\NormalTok{)}
\NormalTok{  ) }\SpecialCharTok{+}
  \FunctionTok{theme\_minimal}\NormalTok{()}

\FunctionTok{print}\NormalTok{(p\_pca)}
\end{Highlighting}
\end{Shaded}

\pandocbounded{\includegraphics[keepaspectratio]{RNA-seq-pipeline_files/figure-latex/plot-1.pdf}}

PCA revealed no clear separation between control and treated samples,
suggesting that the treatment effect is relatively subtle compared to
the overall variability in the dataset.

This may reflect either a mild biological response or inherent
variability between samples.

\subsection{Hierarchical clustering}\label{hierarchical-clustering}

Assess sample similarity.

\begin{Shaded}
\begin{Highlighting}[]
\NormalTok{dist\_mat }\OtherTok{\textless{}{-}} \FunctionTok{dist}\NormalTok{(}\FunctionTok{t}\NormalTok{(df\_log))}
\NormalTok{hc }\OtherTok{\textless{}{-}} \FunctionTok{hclust}\NormalTok{(dist\_mat)}

\FunctionTok{plot}\NormalTok{(hc, }\AttributeTok{main =} \StringTok{"Hierarchical clustering of samples"}\NormalTok{)}
\end{Highlighting}
\end{Shaded}

\pandocbounded{\includegraphics[keepaspectratio]{RNA-seq-pipeline_files/figure-latex/cluster-1.pdf}}
Hierarchical clustering did not fully separate samples by condition,
further supporting the observation that the treatment effect is modest.

Some samples cluster independently of their group, indicating
variability within conditions.

\subsection{Heatmap}\label{heatmap}

Top 10 most highly expressed genes

\begin{Shaded}
\begin{Highlighting}[]
\NormalTok{top\_genes }\OtherTok{\textless{}{-}} \FunctionTok{head}\NormalTok{(}\FunctionTok{order}\NormalTok{(}\FunctionTok{rowMeans}\NormalTok{(df\_log), }\AttributeTok{decreasing =} \ConstantTok{TRUE}\NormalTok{), }\DecValTok{10}\NormalTok{)}

\FunctionTok{pheatmap}\NormalTok{(df\_log[top\_genes, ],}
         \AttributeTok{scale =} \StringTok{"row"}\NormalTok{,}
         \AttributeTok{main =} \StringTok{"Top 10 most expressed genes"}\NormalTok{)}
\end{Highlighting}
\end{Shaded}

\pandocbounded{\includegraphics[keepaspectratio]{RNA-seq-pipeline_files/figure-latex/heatmap-10-1.pdf}}
The heatmap of highly expressed genes shows no strong clustering by
condition, suggesting that the most abundant genes are not strongly
affected by the treatment.

This is consistent with the PCA and clustering results.

\section{4. Differential analysis}\label{differential-analysis}

Organize metadata

\begin{Shaded}
\begin{Highlighting}[]
\NormalTok{coldata }\OtherTok{\textless{}{-}} \FunctionTok{DataFrame}\NormalTok{(}
  \AttributeTok{treatment =} \FunctionTok{factor}\NormalTok{(metadata\_filtered}\SpecialCharTok{$}\NormalTok{TRAT),}
  \AttributeTok{ID =}\NormalTok{ metadata\_filtered}\SpecialCharTok{$}\NormalTok{Identificação,}
  \AttributeTok{row.names =}\NormalTok{ metadata\_filtered}\SpecialCharTok{$}\NormalTok{ID}
\NormalTok{)}

\CommentTok{\#Check order}
\ControlFlowTok{if}\NormalTok{ (}\SpecialCharTok{!}\FunctionTok{identical}\NormalTok{(}\FunctionTok{colnames}\NormalTok{(df), }\FunctionTok{rownames}\NormalTok{(coldata))) \{}
  \FunctionTok{cat}\NormalTok{(}\StringTok{"⚠️Sampless...}\SpecialCharTok{\textbackslash{}n}\StringTok{"}\NormalTok{)}
\NormalTok{  coldata }\OtherTok{\textless{}{-}}\NormalTok{ coldata[}\FunctionTok{colnames}\NormalTok{(df\_filtered), ]}
\NormalTok{\}}
\end{Highlighting}
\end{Shaded}

DESeq2

\begin{Shaded}
\begin{Highlighting}[]
\NormalTok{dds }\OtherTok{\textless{}{-}} \FunctionTok{DESeqDataSetFromMatrix}\NormalTok{(}
  \AttributeTok{countData =}\NormalTok{ df\_filtered,}
  \AttributeTok{colData =}\NormalTok{ coldata, }
  \AttributeTok{design =} \SpecialCharTok{\textasciitilde{}}\NormalTok{ treatment}
\NormalTok{)}

\NormalTok{dds }\OtherTok{\textless{}{-}} \FunctionTok{DESeq}\NormalTok{(dds)}
\end{Highlighting}
\end{Shaded}

Results

\begin{Shaded}
\begin{Highlighting}[]
\NormalTok{res }\OtherTok{\textless{}{-}} \FunctionTok{results}\NormalTok{(dds, }\AttributeTok{contrast =} \FunctionTok{c}\NormalTok{(}\StringTok{"treatment"}\NormalTok{, }\StringTok{"BET"}\NormalTok{, }\StringTok{"CON"}\NormalTok{))}

\NormalTok{res\_summary }\OtherTok{\textless{}{-}} \FunctionTok{data.frame}\NormalTok{(}
  \AttributeTok{Info =} \FunctionTok{c}\NormalTok{(}\StringTok{"p{-}adj \textless{} 0.05"}\NormalTok{, }\StringTok{"p{-}adj \textless{} 0.10"}\NormalTok{, }\StringTok{"p{-}adj \textless{} 0.20"}\NormalTok{),}
  \AttributeTok{Value =} \FunctionTok{c}\NormalTok{(}
    \FunctionTok{sum}\NormalTok{(res}\SpecialCharTok{$}\NormalTok{padj }\SpecialCharTok{\textless{}} \FloatTok{0.05}\NormalTok{, }\AttributeTok{na.rm =} \ConstantTok{TRUE}\NormalTok{),}
    \FunctionTok{sum}\NormalTok{(res}\SpecialCharTok{$}\NormalTok{padj }\SpecialCharTok{\textless{}} \FloatTok{0.10}\NormalTok{, }\AttributeTok{na.rm =} \ConstantTok{TRUE}\NormalTok{),}
    \FunctionTok{sum}\NormalTok{(res}\SpecialCharTok{$}\NormalTok{padj }\SpecialCharTok{\textless{}} \FloatTok{0.20}\NormalTok{, }\AttributeTok{na.rm =} \ConstantTok{TRUE}\NormalTok{)}
\NormalTok{  )}
\NormalTok{)}

\FunctionTok{kable}\NormalTok{(res\_summary) }\SpecialCharTok{\%\textgreater{}\%} \FunctionTok{kable\_styling}\NormalTok{()}
\end{Highlighting}
\end{Shaded}

\begin{longtable}[t]{lr}
\toprule
Info & Value\\
\midrule
p-adj < 0.05 & 0\\
p-adj < 0.10 & 0\\
p-adj < 0.20 & 0\\
\bottomrule
\end{longtable}

\end{document}
